\documentclass[11pt,a4paper]{article}
% Préambule commun aux documents de PythonIPM/documentation.
% Compilation : tectonic <fichier>.tex
% tectonic compile avec XeTeX, qui lit l'UTF-8 nativement : surtout PAS inputenc ni fontenc,
% faits pour pdfLaTeX. Ce sont eux qui décodaient mal les caractères multi-octets — « cœur »
% ressortait en « cur », et les tirets cadratins en caractères parasites.
\usepackage{lmodern}
\usepackage[french]{babel}
\usepackage{microtype}
\usepackage[a4paper,margin=2.4cm]{geometry}
\usepackage{amsmath,amssymb}
\usepackage{booktabs}
\usepackage{longtable}
\usepackage{array}
\usepackage{ragged2e}
\usepackage{xcolor}
\usepackage{tcolorbox}
\usepackage{enumitem}
\usepackage{fancyhdr}
\usepackage[hidelinks]{hyperref}

\definecolor{encadre}{HTML}{F4F6F8}
\definecolor{bordure}{HTML}{4A6FA5}
\definecolor{alerte}{HTML}{B03A2E}

% Encadré « à retenir » : la lecture non technique du passage qui précède
\newtcolorbox{retenir}[1][]{
  colback=encadre, colframe=bordure, boxrule=0.6pt, arc=2pt,
  left=8pt, right=8pt, top=6pt, bottom=6pt,
  fonttitle=\bfseries, title={#1}
}

% Encadré d'avertissement : un point ouvert, un écart assumé
\newtcolorbox{attention}[1][]{
  colback=white, colframe=alerte, boxrule=0.6pt, arc=2pt,
  left=8pt, right=8pt, top=6pt, bottom=6pt,
  % le bandeau de titre a déjà la couleur d'alerte en fond : un titre `coltitle=alerte`
  % écrirait du rouge sur du rouge. On garde le blanc par défaut, comme l'encadré « retenir ».
  fonttitle=\bfseries, title={#1}
}

\newcommand{\col}[1]{\texttt{#1}}
% \ensuremath et non \textsubscript : \ipm est employé aussi bien dans le texte que dans les
% formules, et « M\textsubscript{0} » placé en mode mathématique est une erreur — silencieuse en
% PDF, mais qui fait échouer la conversion vers Word.
\newcommand{\ipm}{\ensuremath{M_{0}}}

\setlist[itemize]{itemsep=2pt, topsep=4pt}
\setlist[description]{style=nextline, leftmargin=1.2cm, itemsep=5pt}

% La source du document apparaît en tête de chaque page. Les documents RGPH la redéfinissent
% par \renewcommand{\source}{...} AVANT \begin{document}.
\newcommand{\source}{EHCVM 2021}

\pagestyle{fancy}
\fancyhf{}
\fancyhead[L]{\small\itshape IPM Côte d'Ivoire --- \source}
\fancyhead[R]{\small\thepage}
\renewcommand{\headrulewidth}{0.3pt}

\renewcommand{\arraystretch}{1.15}

\renewcommand{\source}{PND 2026-2030 / IPM EHCVM 2021}

\title{\bfseries L'IPM calculé sur l'EHCVM 2021 au service du PND 2026-2030\\[4pt]
\large Analyse d'alignement}
\author{Pipeline \col{PythonIPM}}
\date{Septembre 2026}

\begin{document}
\maketitle

\begin{abstract}
\noindent Ce document confronte l'indice de pauvreté multidimensionnelle calculé sur l'EHCVM 2021
au Plan National de Développement 2026-2030. Il ne juge ni le plan ni l'indice : il établit, cible
par cible, ce que l'un permet de mesurer de l'autre, ce qu'il ne permet pas, et l'usage auquel il
se prête le mieux dans le dispositif de suivi du PND. Un document jumeau traite de l'IPM calculé
sur le RGPH 2021 ; un troisième les compare et conclut.
\end{abstract}

\begin{retenir}[La conclusion en cinq lignes]
\begin{itemize}
  \item L'IPM EHCVM couvre \textbf{onze des vingt-deux cibles} du PND qu'un indice de pauvreté
        peut mesurer --- la couverture la plus large des deux sources.
  \item Il est \textbf{le seul} à mesurer la Couverture Maladie Universelle, la sécurité
        nutritionnelle, l'accès effectif aux soins et la protection sociale : quatre cibles
        chiffrées du Pilier 4 auxquelles le recensement ne donne aucune prise.
  \item Il est \textbf{le seul} à porter la pauvreté monétaire et la pauvreté multidimensionnelle
        dans la même base, ce dont l'objectif « baisser significativement le taux de pauvreté » a
        besoin.
  \item Il est \textbf{répété dans le temps} : c'est la seule source capable de dire, en 2028 puis
        en 2030, si le PND a produit ses effets.
  \item Il ne sert en revanche \textbf{pas le ciblage territorial} du Pilier 5 : la moitié de la
        population vit dans une sous-préfecture dont l'IPM n'est pas publiable.
\end{itemize}
\end{retenir}

\tableofcontents
\bigskip

\section{Ce que le PND demande à un indice de pauvreté}

Le PND 2026-2030 vise à « Bâtir une Grande Nation Stable, Ambitieuse et Solidaire ». Parmi ses
objectifs généraux figure de \textbf{« baisser significativement le taux de pauvreté »} --- et
c'est le seul des objectifs majeurs qui ne soit assorti d'aucune valeur cible à 2030, là où le
revenu intermédiaire supérieur, les trois millions d'emplois décents, l'Indice de Capital Humain à
0,7 et les IDE à 3,3~\% du PIB le sont tous.

\begin{retenir}[Une cible manquante, et l'occasion qu'elle ouvre]
Le PND documente la trajectoire de la pauvreté monétaire --- 55,4~\% en 2011, 44,4~\% en 2015,
39,4~\% en 2018, 37,5~\% en 2021 --- mais ne fixe pas de valeur à atteindre en 2030. Un indice de
pauvreté multidimensionnelle est précisément l'instrument qui permettrait de la poser, et de la
suivre sur une grandeur qui ne dépend pas des seuls prix.
\end{retenir}

Le plan formule par ailleurs, pilier par pilier, des cibles chiffrées assorties d'une valeur de
départ et d'une valeur 2030. Ce sont elles qui donnent la mesure de l'alignement : un indice est
utile au PND dans l'exacte mesure où ses indicateurs recoupent ces cibles.

\section{La correspondance, cible par cible}
\label{sec:correspondance}

Les seize indicateurs de l'IPM EHCVM se rapportent aux cibles du PND comme suit. La colonne
« Aujourd'hui » donne la part de la population vivant dans un ménage privé sur l'indicateur.

\subsection{Pilier 4 --- Capital humain, compétences et emplois décents}

\begin{center}
\small
\begin{tabular}{@{}>{\RaggedRight}p{5.3cm}>{\RaggedRight}p{4.6cm}rc@{}}
\toprule
\textbf{Cible PND à 2030} & \textbf{Indicateur IPM} & \textbf{Auj.} & \textbf{Prise} \\
\midrule
Analphabétisme adultes 51,1 $\to$ 45~\% & Alphabétisation & 62,9~\% & directe \\
Analphabétisme 15-24 ans 37,6 $\to$ 20~\% & Alphabétisation & 62,9~\% & directe \\
Achèvement 1\textsuperscript{er} cycle secondaire
  81,14 $\to$ 89~\% & Année de scolarité & 62,9~\% & directe \\
Personnes à moins de 5 km d'un centre de santé 82 $\to$ 100~\% &
  Renoncement aux soins, motif d'indisponibilité de l'offre & 9,0~\% & directe \\
Enrôlés à la CMU 20 $\to$ 30 millions & Assurance maladie & 90,2~\% & directe \\
Couverture par la protection sociale 27 $\to$ 50~\% & Assurance maladie & 90,2~\% & directe \\
Chômage et main-d'œuvre potentielle (SU3) 12,9 $\to$ 8,9~\% & Chômage & 6,5~\% & directe \\
Retard de croissance des moins de 5 ans 21,4 $\to$ 18~\% &
  Insécurité alimentaire (FIES) & 40,8~\% & indirecte \\
Espérance de vie 62,3 $\to$ 65 ans & --- & --- & aucune \\
Indice de l'écart genre 0,655 $\to$ 0,7 &
  désagrégation par sexe du chef de ménage & --- & partielle \\
827 000 ménages bénéficiaires des filets sociaux &
  identification des ménages pauvres & --- & partielle \\
Préscolaire 11 $\to$ 25~\% ; EFTP 6,2 $\to$ 15~\% ; IPS primaire & --- & --- & aucune \\
\bottomrule
\end{tabular}
\end{center}

\subsection{Pilier 5 --- Infrastructures, pôles régionaux et transition écologique}

\begin{center}
\small
\begin{tabular}{@{}>{\RaggedRight}p{5.3cm}>{\RaggedRight}p{4.6cm}rc@{}}
\toprule
\textbf{Cible PND à 2030} & \textbf{Indicateur IPM} & \textbf{Auj.} & \textbf{Prise} \\
\midrule
Accès à une source d'eau améliorée 81 $\to$ 90~\% & Eau potable & 20,9~\% & directe \\
Accès à l'assainissement amélioré 37 $\to$ 70~\% & Toilette & 67,2~\% & directe \\
Accès universel à l'électricité d'ici 2030, en priorité
  dans les quartiers précaires & Électricité & 9,3~\% & directe \\
Habitat décent (défi n\textsuperscript{o}~5 du diagnostic stratégique) &
  Logement, Promiscuité & 17,3~\% ; 12,5~\% & directe \\
Population utilisant internet 40,7 $\to$ 60~\% & Biens d'équipement & 21,7~\% & indirecte \\
Pôles économiques compétitifs, villes secondaires, réduction des
  inégalités régionales & désagrégation territoriale & --- & faible \\
\bottomrule
\end{tabular}
\end{center}

\subsection{Piliers 2, 3 et 6}

\begin{center}
\small
\begin{tabular}{@{}>{\RaggedRight}p{5.3cm}>{\RaggedRight}p{4.6cm}rc@{}}
\toprule
\textbf{Cible ou orientation PND} & \textbf{Indicateur IPM} & \textbf{Auj.} & \textbf{Prise} \\
\midrule
Enregistrement des naissances 56,6 $\to$ 100~\% (P.~6) &
  Déclaration d'état civil & 26,3~\% & directe \\
NNI 27,78 $\to$ 100~\% (P.~6) & Déclaration d'état civil & 26,3~\% & indirecte \\
Modernisation de l'agriculture et productivité (P.~2) ;
  réduction de l'informalité (P.~3) &
  Emploi agricole de subsistance & 40,5~\% & directe \\
Énergie de cuisson propre (transition écologique, P.~5) &
  Énergie de cuisson & 66,3~\% & directe \\
Registre social unifié pour le ciblage (P.~6) & --- & --- & aucune \\
\bottomrule
\end{tabular}
\end{center}

\begin{retenir}[Le décompte]
Sur les vingt-deux cibles et orientations du PND qu'un indice de pauvreté peut éclairer, l'IPM
EHCVM en couvre \textbf{onze de façon directe}, trois de façon indirecte, deux partiellement, et
laisse six hors de portée.
\end{retenir}

\section{Les quatre apports que l'EHCVM est seul à fournir}

\subsection{La Couverture Maladie Universelle et la protection sociale}

Le PND fixe deux cibles chiffrées sur ce terrain : porter les enrôlés à la CMU de 20 à 30 millions,
et la couverture par un dispositif de protection sociale de 27 à 50~\%. L'indicateur
\col{assurance\_maladie} de l'IPM les mesure directement, à l'échelle du ménage.

Le chiffre qu'il donne est d'ailleurs le plus frappant de tout l'indice : \textbf{90,2~\% de la
population vit dans un ménage dont aucun membre n'est couvert}. C'est un ordre de grandeur que la
cible CMU du PND --- 30 millions d'enrôlés sur une population de 29,8 millions --- laisse penser
atteignable, à condition de distinguer l'enrôlement de la couverture effective.

\subsection{La sécurité nutritionnelle}

Le PND fait de la sécurité nutritionnelle un axe stratégique à part entière du Pilier 4, avec une
cible sur le retard de croissance des moins de 5 ans (21,4 $\to$ 18~\%) et une insistance sur les
« 2 000 premiers jours de vie ».

L'EHCVM ne mesure pas l'anthropométrie, mais elle porte l'échelle FIES de la FAO, sur laquelle
l'IPM construit son indicateur d'insécurité alimentaire : \textbf{40,8~\% de la population vit dans
un ménage en insécurité alimentaire modérée ou sévère}. Ce n'est pas le retard de croissance, mais
c'est la mesure du déterminant immédiat que le PND cherche à traiter.

\subsection{L'accès effectif aux soins}

La cible « porter la proportion de personnes vivant à moins de 5 km d'un établissement de santé de
82 à 100~\% » est une cible d'offre. L'indicateur \col{renoncement\_soins} mesure le résultat de
cette offre du point de vue du ménage : un membre a eu un problème de santé, n'a pas consulté, et
a cité comme raison le coût ou \textbf{l'indisponibilité du service ou du personnel}.

C'est ce second motif qui intéresse directement le PND : il mesure le déficit d'offre là où il est
subi, et non la distance théorique à un établissement.

\subsection{Le genre et la pauvreté monétaire}

Deux capacités de désagrégation et de croisement complètent le tableau.

Le \textbf{sexe du chef de ménage} est disponible dans l'EHCVM, ce qui permet de servir la cible
« indice de l'écart genre 0,655 $\to$ 0,7 » sur le terrain de la pauvreté multidimensionnelle.

Surtout, l'EHCVM porte dans la même base la \textbf{consommation des ménages} : c'est elle qui
produit le taux de pauvreté monétaire cité par le PND. Les deux pauvretés peuvent donc être
croisées ménage par ménage.

Le tableau ci-dessous porte sur les 9\,800 ménages pour lesquels l'indicateur de bien-être est
renseigné, soit 18,7 millions de personnes.

\begin{center}
\begin{tabular}{@{}lrr@{}}
\toprule
& \textbf{Pauvre au sens de l'IPM} & \textbf{Non pauvre IPM} \\
\midrule
Pauvre monétaire & 41,7~\% & 6,7~\% \\
Non pauvre monétaire & 28,2~\% & 23,4~\% \\
\bottomrule
\end{tabular}
\end{center}

Deux enseignements, dont un seul est solide.

Le premier l'est : \textbf{86,1~\% des pauvres monétaires sont aussi pauvres au sens
multidimensionnel}. C'est un rapport calculé à l'intérieur du champ, donc insensible à la façon
dont on traite les ménages sans indicateur de bien-être --- il vaut 86,1~\% quelle que soit la
convention retenue. La pauvreté monétaire est presque entièrement incluse dans la pauvreté
multidimensionnelle.

Le second ne l'est pas au même degré, mais sa direction ne fait pas de doute : une part importante
de la population est \textbf{privée sans être pauvre monétairement} --- 28,2 points ici, 31,7
points si l'on raisonne sur l'ensemble de l'échantillon. Une politique qui ne suivrait que le
revenu manquerait cette population, que le PND entend pourtant servir.

\begin{attention}[Pourquoi les niveaux ne doivent pas être publiés tels quels]
L'indicateur de bien-être (\col{pcexp}) et le seuil national (\col{zref}) ne sont renseignés que
pour \textbf{9\,800 ménages sur 12\,965}, soit 75,6~\% de l'échantillon et 62,9~\% de la population
pondérée : le module de consommation n'est pas exploitable pour un quart des ménages.

Aucune des deux lectures possibles ne retombe sur le chiffre publié :

\begin{center}
\begin{tabular}{@{}lr@{}}
\toprule
\textbf{Lecture} & \textbf{Taux monétaire} \\
\midrule
Manquants comptés non pauvres, sur toute la population & 30,5~\% \\
Champ restreint aux 9\,800 ménages renseignés & 48,5~\% \\
\textbf{Publié, et repris par le PND} & \textbf{37,5~\%} \\
\bottomrule
\end{tabular}
\end{center}

Et le sous-champ n'est pas neutre : l'IPM y vaut 69,9~\% contre 58,0~\% sur l'échantillon complet.
\textbf{Les ménages disposant d'un indicateur de bien-être sont plus pauvres que la moyenne.}
Reproduire le 37,5~\% demanderait la pondération ou l'imputation que l'INS applique à ce
sous-champ, qui ne figurent pas dans la base.

Les déflateurs spatial et temporel présents dans la base (\col{def\_spa}, \col{def\_temp}) ne sont
pas en cause : quelle que soit la façon de les appliquer, le taux reste entre 30,3 et 30,9~\%.

\textbf{À retenir} : le taux de recouvrement de 86,1~\% est publiable, les quatre cases du tableau
ne le sont pas avant que la méthode officielle de calcul du taux monétaire ait été reprise.
\end{attention}

\section{Ce que l'IPM EHCVM ne peut pas faire pour le PND}

\subsection{Le ciblage territorial du Pilier 5}

Le Pilier 5 repose sur une « planification spatiale des investissements, fondée sur une
cartographie fine des potentialités économiques régionales », sur des Pôles Économiques Compétitifs
et sur le développement des villes secondaires. Le Pilier 6 prévoit un « registre social unifié
visant à renforcer l'inclusion et le ciblage des politiques publiques ». Le Pilier 4 fixe un
objectif de 827 000 ménages bénéficiaires de transferts monétaires.

Ces trois orientations demandent une information à la maille du territoire et, pour la dernière, à
la maille du ménage. L'EHCVM est un échantillon de 12 965 ménages : elle ne peut satisfaire ni
l'une ni l'autre.

\begin{attention}[La moitié de la population hors de portée d'un chiffre publiable]
Sur les 442 sous-préfectures et communes que l'EHCVM permet de distinguer, \textbf{92 --- soit
47,9~\% de la population} --- ont un IPM dont le coefficient de variation dépasse 20~\% : le
chiffre existe, mais il n'est pas publiable.

À la maille du département, 15 des 108 sont dans ce cas, couvrant 7,1~\% de la population.

Autrement dit : l'EHCVM donne un IPM solide pour le pays, pour les deux milieux et pour les 33
régions, mais elle ne peut pas fonder une allocation d'investissements à la maille où le Pilier 5
entend la conduire.
\end{attention}

\subsection{Les cibles hors champ}

Six cibles du PND restent sans prise : l'espérance de vie, le taux net de scolarisation au
préscolaire, la part de l'EFTP dans le secondaire, l'indice de parité filles/garçons dans le
primaire, le NNI et le registre social unifié. Les quatre premières relèvent de statistiques
sectorielles ; les deux dernières, d'un système d'information administratif.

\section{L'usage auquel l'IPM EHCVM se prête}

\subsection{Le tableau de bord du cadre de résultats}

C'est son emploi naturel. Onze cibles du PND, réparties sur quatre piliers, peuvent être suivies
sur un indicateur unique, décomposable, dont la contribution de chaque dimension est explicite.

\begin{center}
\begin{tabular}{@{}lrr@{}}
\toprule
\textbf{Dimension} & \textbf{Poids nominal} & \textbf{Contribution à \ipm} \\
\midrule
Éducation & 25~\% & 32,3~\% \\
Santé & 25~\% & 27,9~\% \\
Conditions de vie & 25~\% & 21,3~\% \\
Emploi & 25~\% & 18,5~\% \\
\bottomrule
\end{tabular}
\end{center}

\noindent Un gouvernement qui veut savoir quel levier déplacerait le plus l'indice dispose là de la
réponse, et elle est directement transposable en priorités de pilier.

\subsection{Le suivi dans le temps --- l'argument décisif}

Le PND court sur cinq ans et prévoit un cadre de gouvernance chargé d'en mesurer les progrès. Or
l'EHCVM est une \textbf{enquête répétée} : elle existe pour 2018 et 2021, et une nouvelle vague
interviendra dans l'intervalle du plan.

C'est la seule des deux sources qui permette de dire, à mi-parcours puis en 2030, si les cibles ont
été atteintes. Le recensement, décennal, ne le permettra pas avant le plan suivant.

\begin{retenir}[Trois emplois immédiats]
\begin{enumerate}
  \item \textbf{Poser la cible manquante.} L'objectif « baisser significativement le taux de
        pauvreté » peut être chiffré sur \ipm{} = 0,2863 (2021), avec une valeur 2030 explicite.
  \item \textbf{Alimenter le cadre de résultats} sur les onze cibles à prise directe, en une seule
        série cohérente plutôt qu'en onze statistiques sectorielles.
  \item \textbf{Arbitrer entre piliers} au vu des contributions : la dimension Santé pèse 28,4~\%
        de l'indice, la dimension Emploi 18,2~\%.
\end{enumerate}
\end{retenir}

\section{Limites de cette analyse}

\begin{enumerate}
  \item \textbf{Les correspondances ne sont pas des identités.} « Alphabétisation » au sens de
        l'IPM est un indicateur de ménage --- un membre de 17-49 ans ne sait pas lire ou écrire le
        français --- tandis que la cible PND porte sur un taux individuel. Les deux évoluent dans
        le même sens sans être le même nombre.
  \item \textbf{La dimension Santé de l'IPM EHCVM est portée à 93,4~\% des ménages par l'assurance
        maladie}, ce qui en fait un indicateur peu discriminant (voir la \emph{Note sur la
        dimension santé}). Son poids dans l'indice doit être lu avec cette réserve.
  \item \textbf{Le croisement monétaire / multidimensionnel ne reproduit pas le taux de pauvreté
        publié}, parce que l'indicateur de bien-être manque pour un quart des ménages et que ce
        quart n'est pas tiré au hasard. Seul le taux de recouvrement entre les deux pauvretés
        (86,1~\%) est robuste à ce défaut.
  \item \textbf{Le PND ne contient pas de cible d'IPM.} L'alignement établi ici est un alignement
        de \emph{contenu}, pas une correspondance institutionnelle : il resterait à faire adopter
        l'indice dans le cadre de gouvernance du plan.
\end{enumerate}

\end{document}
